\section{Method}
\subsection{Problem Formulation}
Let $i=1,\dots,N$ index award notices in the analysis frame and let $X_i\in\{0,1\}$ denote AI-relatedness as defined by a fixed codebook. $X_i$ is observed only for a labelled subsample $\mathcal{L}$ drawn with known inclusion probabilities $\pi_i$; for every notice, a method $m$ produces a prediction $\hat{X}^m_i$ or a score. We target (a) the precision, recall, and F1 of each method relative to $X$, (b) the prevalence $\theta=N^{-1}\sum_i X_i$, and (c) the coefficient $\beta$ on $X_i$ in a model of single bidding $Y_i$ with controls $W_i$. The reference annotation defines the construct operationally; it is not treated as expert ground truth.

\subsection{Construct and Reference Annotation}
A notice is AI-related if the procured object includes the development, supply, licensing, operation, or consulting and training on a system whose core functionality relies on AI techniques (machine learning, deep learning, natural language processing or LLMs, chatbots, computer vision, speech recognition or synthesis, or predictive models learned from data). Generic software, workflow automation without AI, dashboards without learning components, and hardware are coded as not AI-related. This definition follows the task-based view of AI used in vacancy studies \cite{acemoglu2022} and the procurement-risk framing of AI systems \cite{hickok2022}. Each sampled notice was annotated by a primary LLM (Claude Haiku 5.5) and a secondary LLM (Gemini 2.5 Flash-Lite) with the same zero-shot prompt (temperature 0, JSON output); the primary label is the reference, and agreement is reported as Cohen's $\kappa$. LLM annotation is used because it has shown accuracy comparable to or above crowd workers on several tasks \cite{gilardi2023}, while its task dependence \cite{ziems2024} motivates the two-annotator check and the sensitivity analyses in Sect.~\ref{sec:robust}.

\subsection{Measurement Methods}
\textbf{M1} applies an English AI dictionary (22 patterns, e.g., \textit{artificial intelligence}, \textit{machine learning}, \textit{chatbot}, \textit{LLM}) to every notice, mirroring common practice. \textbf{M2} adds core equivalents of \textit{artificial intelligence}, \textit{machine learning}, \textit{neural network}, \textit{language model}, and \textit{chatbot} in the EU languages and the abbreviations KI and IA. \textbf{M3} is zero-shot NLI with mDeBERTa-v3-base-mnli-xnli \cite{he2023,yin2019} and the hypothesis ``This procurement is about artificial intelligence'' (entailment probability $\ge 0.5$). \textbf{M4} is the maximum cosine similarity between a notice and five English prototype queries in the multilingual-e5-small space \cite{wang2024me5}, thresholded within training folds. \textbf{M5}, the proposed measurement, is an $L_2$-regularised logistic regression on multilingual-e5-small embeddings trained on the reference labels with weights $1/\pi_i$. \textbf{M6} replaces the embeddings with character 3--5-gram TF-IDF features, a strong classical baseline for short multilingual text. M5 and M6 represent the supervised route recommended when labels are available \cite{gentzkow2019}; M3 and M4 require no labels. All label-trained components are cross-fitted: predictions for labelled notices come from five folds that exclude them, which prevents optimistic evaluation and is required for DSL \cite{egami2023}.

\subsection{Sampling Design and Design-Based Estimation}
Because AI notices are rare, a simple random sample would contain few positives. The frame was partitioned into four strata: A (M1 hit), B (M2 hit but not M1), C (neither, but M4 similarity in the top 3\%), and D (remaining notices), and simple random samples were drawn within strata (seed 2027). Accuracy metrics weight each labelled notice by $1/\pi_i$, which makes them estimates for the whole frame. For prevalence and regression, DSL replaces the unobserved $X_i$ with the pseudo-outcome
\begin{equation}
\tilde{X}_i=\hat{g}(T_i)+\frac{R_i}{\pi_i}\bigl(X_i-\hat{g}(T_i)\bigr),
\end{equation}
where $\hat{g}$ is the cross-fitted M5 probability, $T_i$ the text, and $R_i$ the labelling indicator. Because $E[\tilde{X}_i\mid T_i]=E[X_i\mid T_i]$ under the known design, $\hat\theta_{\mathrm{DSL}}=N^{-1}\sum_i\tilde{X}_i$ is unbiased whatever the classifier error. For the linear probability model $Y_i=\beta X_i+W_i^\top\gamma+\varepsilon_i$, we substitute $\tilde{X}_i$ for $X_i$ and for $X_i^2$ in the ordinary least squares moment conditions (valid because $X_i$ is binary) and use sandwich standard errors. We use a linear probability model because the specification contains over 100 fixed effects, for which logit estimates are less stable, and report logit plug-in estimates as a robustness check.
